每个测试配置，即\emph{配置单元（cell）}，固定数据集、索引和工作负载。组件比较改变待研究的机制，并在适用情况下固定训练索引、查询、probe 数、精排预算和评分表示；跨系统比较则使用各系统自己的参数网格（\cref{sec:7-2}）。\Cref{tab:3-design} 将研究问题对应到这些对照和测量指标。成对测试衡量执行改变带来的效果，而缓存行为等硬件原因需要单独剖析。

以 $Q$ 表示每秒查询数（QPS），成对加速比为 $Q_{\rm treatment}/Q_{\rm control}$，大于一时表示处理方案更快。对 $J$ 个配置使用比值的几何平均，各配置权重相同：
\[
 G=\exp\!\left(\frac1J\sum_{j=1}^{J}
 \log\frac{Q_{{\rm treatment},j}}{Q_{{\rm control},j}}\right).
\]

\par\medskip\noindent\begin{minipage}{\linewidth}
\centering\captionsetup{hypcap=false}
\captionof{table}[研究问题与实验对照]{研究问题、受控比较与检查。构建、执行、预算选择和跨系统搜索需要不同的对照。}
\label{tab:3-design}
\footnotesize\setlength{\tabcolsep}{3pt}
\begin{tabularx}{\linewidth}{@{}>{\raggedright\arraybackslash}p{.19\linewidth}*{3}{>{\raggedright\arraybackslash}X}@{}}
\toprule
研究问题 & 比较 & 共同条件 & 指标与检查 \\
\midrule
RQ1：GPU 索引构建与容量 & 冷构建与训练模型复用；压缩与原始存储
& $N,d$、表示、路由及构建计时边界
& 训练／编码时间及常驻／峰值显存；实测与模型估算的显存对比 \\
RQ2：面向表示结构的执行 & 查表分组、打包读取、协作精排、分组扫描、共享阈值与 probe 选择
& 相同索引、负载与预算；共同的评分定义
& 主机调用 QPS、阶段时间，以及按对应并列规则核对的分数和返回 ID 差异 \\
RQ3：工作负载级精排校准 & 无标签抽样／全池校准；固定 64、100、200；完整预算网格搜索及使用 GT 的诊断选择
& 相同索引与路由；分开的校准和评估查询
& 测试查询相对最佳已测预算的召回损失、所选预算 QPS 与校准成本；检验质量与回本 \\
RQ4：跨系统搜索性能 & 不同 batch 大小下分别调优的 JHQ、CAGRA、IVF-PQ 与 IVF-RaBitQ；暴力搜索参照
& 相同导出向量、GT、$k$、GPU 和主机调用边界；显存单独报告
& 匹配召回率 QPS、显存与覆盖范围；记录胜负 \\
\bottomrule
\end{tabularx}
\end{minipage}\par\medskip
